\section{Einleitung}
\subsection{Vorstellung des Unternehmens}
Die Brau Union Österreich AG (BUÖ) ist das größte Brauereiunternehmen Österreichs und entstand 1997 durch den Zusammenschluss der Brau AG in Linz und der Steirerbrau AG in Graz. Die Unternehmensgeschichte reicht teilweise bis ins 15. Jahrhundert zurück. Mit einem Marktanteil von über 50 \% nimmt die Brau Union eine führende Stellung auf dem österreichischen Biermarkt ein. 

Im Jahr 2001 erzielte das Unternehmen einen Umsatz von rund 600 Millionen Euro und einen Getränkeabsatz von etwa 5,7 Millionen Hektolitern. Zudem beschäftigte die Brau Union rund 2.500 Mitarbeiter und verfügte über einen eigenen Fuhrpark von etwa 400 Lkw. 

Die BUÖ ist Teil der Österreichischen Brau-Beteiligungs-AG (BBAG), zu der neben österreichischen Brauereien auch weitere Getränkeunternehmen gehören. Ein zentraler Bestandteil der Unternehmensphilosophie ist die Mehrmarkenstrategie, mit der unterschiedliche Konsumentenbedürfnisse über verschiedene Marken und Preissegmente abgedeckt werden. Der Vertrieb erfolgt hauptsächlich über den Lebensmittelhandel sowie die Gastronomie. 

Das Controlling hat im Unternehmen eine lange Tradition und unterstützt das Management bei der Planung und Steuerung der Unternehmensziele. Dabei arbeiten Manager und Controller eng zusammen, um sowohl strategische als auch operative Ziele zu erreichen.
\subsection{Strategische Ausgangslage der BUÖ}

Die BUÖ verfolgt eine konsequente Mehrmarkenstrategie, um unterschiedliche Konsumentenbedürfnisse in verschiedenen Preis- und Image-Segmenten abzudecken. Der Vertrieb erfolgt über die zentralen Kanäle Lebensmittelhandel und Gastronomie. Mit einem Marktanteil von über 50 \% nimmt das Unternehmen eine dominante Stellung im österreichischen Biermarkt ein.
Das Controlling ist dezentral organisiert und auf einem hohen Entwicklungsniveau. Es unterstützt das Management sowohl in operativen als auch in strategischen Fragestellungen. Die wertorientierte Steuerung basiert insbesondere auf Kennzahlen wie Return on Investment, Unternehmensergebnis und Free Cashflow. Vor diesem Hintergrund ist eine Übersetzung dieser strategischen Zielgrößen in operative Werttreiber erforderlich.

\subsection{Bedeutung der Balanced Scorecard für die BUÖ}

Die Balanced Scorecard ist ein Managementinstrument, das Unternehmen dabei unterstützt, ihre strategischen Ziele in konkrete Maßnahmen und Kennzahlen zu übersetzen. Neben finanziellen Kennzahlen berücksichtigt sie auch weitere Perspektiven wie Kunden, interne Prozesse sowie Lernen und Entwicklung (vgl. \cite{gadatsch_it-controlling_2026} S. 73-74).

Für die Brau Union Österreich AG ist dies besonders relevant, da sich der Biermarkt zunehmend verändert. In vielen europäischen Ländern ist ein Rückgang des Bierkonsums zu beobachten (vgl. \cite{noauthor_bierabsatz_nodate}).

Die Balanced Scorecard ermöglicht es der Brau Union, strategische Ziele systematisch umzusetzen und sowohl finanzielle als auch nicht-finanzielle Erfolgsfaktoren zu berücksichtigen. Dadurch wird eine ganzheitliche Steuerung des Unternehmens unterstützt und die langfristige Wettbewerbsfähigkeit gesichert.